% ##############################################################################
\section{Background, problem statement, and threat model}
\label{sec:background}

This section fixes the notation inherited from the white paper
\cite{noesis_whitepaper}, states how the prototype changes it, and formulates
the probing problem.

% ==============================================================================
\subsection{Contracts and fulfillment monitoring in the white paper}
\label{sec:background_whitepaper}

The white paper defines an \emph{intelligence contract} as the tuple
$$
  \kappa = (\ntasksattr,\ \clevelattr,\ \lmaxattr,\ \rminattr,\ \priceattr),
$$
where $\clevelattr$ is a required capability tier in a totally ordered set
$\calC$, $\lmaxattr$ is the maximum latency per task, and $\rminattr \in [0,1]$
is the minimum fraction of tasks that must meet both $\clevelattr$ and
$\lmaxattr$. For a request $x$ served under $\kappa$, the monitoring layer
estimates the delivered tier $\hat c(x)$ and the measured latency
$\hat\ell(x)$, and defines the success indicator
\begin{equation}
  \label{eq:wp_success}
  s(x) = \mathbb{1}\!\left[\hat c(x) \succeq \clevelattr(\kappa)
    \;\wedge\; \hat\ell(x) \leq \lmaxattr(\kappa)\right].
\end{equation}
The measured reliability over the request set $W$ served during the
fulfillment window is the sample mean $\hat r(\kappa, W) = |W|^{-1}\sum_{x\in W}
s(x)$. To separate under-delivery from sampling noise, the white paper treats
monitoring as a one-sided hypothesis test with the null hypothesis that the
seller is compliant, and flags $\kappa$ when a lower confidence bound
$\hat r_{\text{lower}}(\kappa, W)$ falls below $\rminattr(\kappa)$. The estimator
of $\hat c(x)$ is a pluggable component, and the white paper names a verifier
model, a benchmark probe, and agreement with a stronger reference model as
candidates.

% ==============================================================================
\subsection{Refinements in the prototype}
\label{sec:background_refinements}

The prototype keeps the architecture of the white paper and changes three
details that matter for the canary subsystem.

\begin{enumerate}
  \item \textbf{Capability is measured by known-answer probes.} The estimator
    $\hat c(x)$ is replaced by a binary grade on a canary question. The
    prototype does not map the grade to a tier in $\calC$.
  \item \textbf{The joint indicator is factored.} The indicator
    \eqref{eq:wp_success} is split into a latency metric over buyer requests
    and a quality metric over canary requests. Let $B(\kappa,W)$ denote the
    buyer requests and $C(\kappa,W)$ the canary requests of a window. The
    latency success of $x \in B$ is $s_\ell(x) = \mathbb{1}[\hat\ell(x) \leq
    \lmaxattr]$ for a request that the seller answered, and $s_\ell(x) = 0$ for
    a seller timeout or upstream error. The quality success of $x \in C$ is
    $s_q(x) = \grade(c_x, y_x)$, where $c_x$ is the canary and $y_x$ the
    response, and $s_q(x) = 0$ for a seller timeout or upstream error. The contract carries a promised canary accuracy
    $\qminattr \in [0,1]$ in addition to $\rminattr$, which applies to latency.
  \item \textbf{The violation rule uses an upper bound.} A metric is a
    violation only when the one-sided confidence \emph{upper} bound on its true
    rate is below the promise. Section~\ref{sec:verdict_rule} analyzes this
    choice against the lower-bound rule.
\end{enumerate}

Table~\ref{tab:wp_mapping} summarizes the correspondence.

\begin{table}[htbp]
  \centering
  \small
  \begin{tabularx}{\linewidth}{@{}p{0.2\linewidth}XX@{}}
    \toprule
    Concept & White paper & Prototype \\
    \midrule
    Capability estimate & pluggable $\hat c(x)$ (verifier, probe, reference
      agreement) & binary exact-match grade on a canary, behind a
      \texttt{CapabilityMeasurer} interface \\
    Success indicator & one joint indicator $s(x)$ & latency indicator on
      buyer requests, quality indicator on canaries \\
    Promise & $\clevelattr$ and $\rminattr$ & $\rminattr$ for latency,
      $\qminattr$ for canary accuracy \\
    Violation test & lower bound below $\rminattr$ & upper bound below the
      promise \\
    Evidence minimum & not specified & 20 buyer requests and 8 canaries \\
    Attribution & request tagged with $\kappa$ & provider pinned per contract,
      mismatches excluded \\
    \bottomrule
  \end{tabularx}
  \caption{Correspondence between the white paper's monitoring model and the
  prototype's canary-based implementation.}
  \label{tab:wp_mapping}
\end{table}

% ==============================================================================
\subsection{Problem statement}
\label{sec:problem}

\begin{problem}[Probing a seller's delivered quality]
\label{prob:probing}
Given an active contract $\kappa$ with seller $s$, a collection window of
length $W$ seconds that opens at time $t_0$, a target number of probes $N$, a
secret key $k$, and a bank of candidate canaries $\bank$, produce a sequence of
probes $(c_i, t_i)_{i=1}^{N}$ with $c_i \in \bank$ and $t_i \in [t_0, t_0+W]$
such that the properties below hold.
\end{problem}

\begin{description}
  \item[P1 Validity.] The expected answer of every canary is known by
    construction, and grading is a deterministic function of the response.
  \item[P2 Discrimination.] A canary answered correctly by every healthy
    provider is answered incorrectly with non-negligible probability by a
    weaker model, so that correct answers carry information about the
    contracted capability.
  \item[P3 Indistinguishability at the provider boundary.] The request that the
    seller receives contains no field that marks it as a probe.
  \item[P4 Unpredictability.] Without $k$, neither the identity nor the timing
    of a probe is predictable from earlier probes.
  \item[P5 Faithful attribution.] A probe outcome counts against the seller
    only if the seller caused it. Gateway, configuration, grader, and buyer
    errors are not charged to the seller.
  \item[P6 Auditability.] Every probe, response, and grade is persisted, and the
    same inputs and key reproduce the same selection and schedule.
\end{description}

% ==============================================================================
\subsection{Threat model}
\label{sec:threat_model}

The adversary is a seller that wishes to lower its serving cost while keeping
the contract revenue. Table~\ref{tab:threats} lists the behaviors considered
and indicates whether the canary subsystem is intended to address them.

\begin{table}[htbp]
  \centering
  \small
  \begin{tabularx}{\linewidth}{@{}p{0.22\linewidth}XX@{}}
    \toprule
    Seller behavior & Effect on evidence & Handled by canaries? \\
    \midrule
    Serves a weaker model under the contracted name & lower accuracy on
      canaries & Yes, if the substitute scores low on the calibrated bank
      (Section~\ref{sec:detection_power}) \\
    Serves a quantized or otherwise slightly degraded variant & small accuracy
      change & Weakly: power is low when the shortfall from the promise is small
      \\
    Adds latency & buyer-request latency rises & No: canary latency is
      excluded and latency is judged on buyer traffic \\
    Recognizes probes and serves them with a stronger model & canaries
      overstate quality & Not prevented; the prompt distribution is a known
      weakness (Section~\ref{sec:limitations}) \\
    Times out or returns upstream errors & recorded as seller failures & Yes,
      counted as failed quality evidence \\
    \bottomrule
  \end{tabularx}
  \caption{Seller behaviors considered and whether the canary subsystem is
  designed to detect them.}
  \label{tab:threats}
\end{table}

The analysis rests on the following trust assumptions. The gateway, the
database, and the checker are operated by the platform and are trusted. The
key $k$ is held only by the platform. The bank content is \emph{not} assumed
secret in the prototype, because the bank is checked into the research
repository and a deterministic script in the same repository regenerates it; the
secrecy that the design provides concerns which canaries are selected for a
contract and when they are sent. Provider pinning
is assumed to work: the gateway routes each contract to the provider slug of the
seller and records the provider that actually answered. Collusion between
sellers and between a seller and a buyer is out of scope.
